\begin{dang}{CHU KÌ VÀ TẦN SỐ CỦA CON LẮC LÒ XO}
%\begin{itemize}
%\item Dựa vào mối quan hệ thuận nghịch để rút ra biểu thức liên hệ. 
%\item Dấu $\sim$ ứng với tỉ lệ.
%\item Phép thế như dưới đây gọi là phép thế theo tỉ lệ sẽ rút ngắn được thời gian biến đổi.
%\end{itemize}
%\setcounter{paragraph}{0}
%$\bullet$ Sự thay đổi chu kì, tần số theo khối lượng
%\begin{itemize}
%\item Ta có: \begin{align*}
%T_1=\dfrac{1}{f_1}=2\pi \sqrt{\dfrac{m_1}{k}}.\quad T_2=\dfrac{1}{f_2}=2\pi \sqrt{\dfrac{m_2}{k}}.
%\end{align*}
%\item Khi $m=am_1 \pm bm_2$ với $a,\ b$ là hằng số:
%\begin{itemize}
%\item Do $m\sim T^2 \xrightarrow{\text{thế}\  T^2\ \text{tương ứng m}}$ $T^2=aT_1^2\pm bT_2^2$. 
%\item Do $m\sim \dfrac{1}{f^2} \xrightarrow{\text{thế}\  \dfrac{1}{f^2}\ \text{tương ứng m}}$ $\dfrac{1}{f^2}=\dfrac{a}{f_1^2}\pm\dfrac{b}{f_2^2}$. 
%\end{itemize} 
%\end{itemize}
%$\bullet$ Sự thay đổi chu kì, tần số theo độ cứng k
%\begin{itemize}
%\item Ta có: \begin{align*}
%T_1=\dfrac{1}{f_1}=2\pi \sqrt{\dfrac{m}{k_1}};\quad T_2=\dfrac{1}{f_2}=2\pi \sqrt{\dfrac{m}{k_2}}.
%\end{align*}
%\item Khi $k=ak_1 \pm bk_2$ với $a,\ b$ là hằng số:
%\begin{itemize}
%\item Do $k \sim \dfrac{1}{T^2} \xrightarrow{\text{thế}\  \dfrac{1}{T^2}\ \text{tương ứng k}}$ $\dfrac{1}{T^2}=\dfrac{a}{T^2}\pm \dfrac{b}{T^2}$. 
%\item Do $k \sim f^2 \xrightarrow{\text{thế}\ f^2\ \text{tương ứng k}\ \ \ \ \  }$ $f^2=af_1^2\pm bf_2^2$.
%\end{itemize} 
%\end{itemize} 
\end{dang}